\section{Conclusion}
\label{sec:conclusion}

In this work, we investigated semantic correspondence using features extracted from modern visual foundation models. Through extensive experiments
on the SPair-71k benchmark, we analyzed the inherent correspondence capabilities of several pretrained backbones, as well as the impact of light
fine-tuning and improved prediction strategies. Our results demonstrate that foundation models already encode strong semantic alignment cues,
enabling accurate correspondence estimation even in a training-free setting.

Among the evaluated backbones, DINO-based models substantially outperform Segment Anything for semantic correspondence. While DINOv3 achieves the
highest accuracy at moderate thresholds, reaching 75.77\% at PCK@0.1, DINOv2 attains superior performance at stricter thresholds, achieving 60.02\% at
PCK@0.05. These results highlight a clear trade-off between localization precision and robustness across different pretrained representations.
We further showed that selectively fine-tuning the final layers of the backbone using a contrastive objective consistently improves performance
across all thresholds, confirming that limited adaptation of pretrained features can significantly enhance correspondence quality. In addition,
replacing the standard argmax operation with a windowed soft-argmax prediction leads to more robust and precise keypoint localization.

To assess generalization across domains, we additionally evaluated our approach on the PF-Willow dataset.
Despite substantial differences in object categories, scale variation, and annotation protocols compared to SPair-71k,
DINO-based models maintain strong correspondence performance, exceeding 91\% PCK@0.1 without any additional adaptation.
DINOv3 consistently outperforms DINOv2 across most categories, while SAM exhibits a significant performance drop,
confirming the superior transferability of self-supervised transformer representations for correspondence tasks.
These results indicate that the learned representations capture transferable semantic structures that generalize beyond the training distribution of SPair-71k.

Overall, our findings highlight the effectiveness of visual foundation models as a robust and scalable solution for semantic correspondence across
datasets and domains. This study provides practical insights into model selection, fine-tuning strategies, and prediction design, and suggests that
further gains may be achieved by exploring domain adaptation techniques or extending the approach to related tasks such as geometric correspondence
and point tracking in video.

